\section{Estimation and Forecasting}
Estimation is a tool for planning, not a promise. Use points to compare relative size, and use throughput for forecasting. Keep the system honest by measuring what actually finishes.

\subsection{Story points vs throughput}
Story points help with sprint sizing and team learning. Throughput (items completed per sprint) is the strongest predictor of near-term capacity. Use points for in-sprint planning and throughput for external commitments.

Sample JQL filters:
\begin{itemize}
  \item Completed last sprint: \texttt{project = PMH AND sprint in closedSprints() AND status = Done}
  \item Unestimated in current sprint: \texttt{project = PMH AND sprint in openSprints() AND Effort is EMPTY}
\end{itemize}

\subsection{Probabilistic forecasting (no heavy math)}
Use past throughput to create ranges. Example: if the team finishes 8--12 items per sprint, forecast delivery in 2--3 sprints for a 20-item Epic. The goal is a confidence range, not a single date. Update the range each sprint based on real throughput.

\subsection{Case study insert}
Estimates for ST-105 through ST-108 are used to size EPIC-2. Throughput from the last three sprints informs how much import work can fit in R1.
\subsection{Release burndown pitfalls}
Burndowns lie when scope changes. If scope grows, the burndown looks worse even if delivery is steady. Always show scope change alongside burn rate.

Sample JQL filters:
\begin{itemize}
  \item Scope added mid-sprint: \texttt{project = PMH AND sprint in openSprints() AND createdDate >= startOfSprint()}
\end{itemize}









